\chapter[Control Flags]{Control Flags}
\label{cp:control-flags}


\begin{listing}[!htpb]
    \caption{Defining the control flags.}
    \label{listing:control-flags}
    \begin{minted}{matlab}
%% Control Flags
N_bits         = 100;           % Number of bits per realization.
N_realizations = 500;           % Number of realizations.
Pw             = 0.07;          % Pulse width of the bit.
Ts             = 0.01;          % Time sample.
L              = round(Pw/Ts);  % Samples per symbol period (L=7).
A              = 4;             % Amplitude.
N_fft          = 1024;          % FFT size.
Fs             = 1/Ts;          % Sampling frequency.
freq           = (-N_fft/2:N_fft/2-1) * (Fs/N_fft);  % Frequency axis.
center         = floor(N_fft/2) + 1;                  % Center bin.
half_len       = 100;           % Max tau used in correlation_manual.
\end{minted}
\end{listing}


\begin{itemize}
    \item \texttt{N\_bits}: Defines the number of bits generated in each realization of the random process. It controls the length of the bit stream.
    \item \texttt{N\_realizations}: Specifies the number of independent realizations of the random process.
    \item \texttt{Pw}: Represents the pulse width (bit duration) in seconds. It determines the time-domain duration of each bit and directly affects the signal's bandwidth.
    \item \texttt{Ts}: The time sampling interval in seconds.
    \item \texttt{L}: It determines how many discrete samples represent a single bit.
    \item \texttt{A}: Sets the peak amplitude of the transmitted signal.
    \item \texttt{N\_fft}: Specifies the size of the Fast Fourier Transform (FFT).
    \item \texttt{Fs}: The sampling frequency ($1/Ts$). It defines the frequency range of the simulation and must be high enough to satisfy the Nyquist criterion.
    \item \texttt{freq}: The frequency axis vector used for spectral analysis, mapped from $-Fs/2$ to $Fs/2$ to allow for centered Power Spectral Density (PSD) plots.
    \item \texttt{center}: The index of the zero-frequency (DC) component in the FFT output, used to correctly align and shift the spectrum for plotting.
    \item \texttt{half\_len}: Defines the maximum lag ($\tau$) for the manual calculation of the autocorrelation function.
\end{itemize}